% 1990 Q19 质点圆弧路径与变力做功图 (解析几何)
% 坐标轴
\draw[->, line width=0.9pt] (-0.5, 0) -- (4.8, 0) node[below=1pt] {$x$};
\draw[->, line width=0.9pt] (0, -0.5) -- (0, 5.2) node[left=1pt] {$y$};
\node[below left=1pt] at (0,0) {$O$};

% 点 A(1,2) 与 点 B(3,4)
\coordinate (A) at (1, 2);
\coordinate (B) at (3, 4);
\coordinate (C) at (2, 3);

% 直径 AB (虚线)
\draw[dashed, line width=0.9pt] (A) -- (B);

% 以 AB 为直径的半圆弧 (圆心 (2,3), 半径 sqrt(2) ≈ 1.414, 角度从 -135° 到 45°)
\draw[line width=1.1pt] (A) arc[start angle=-135, end angle=45, radius=1.4142];

% 标记顶点
\node[left=2pt] at (A) {$A(1,2)$};
\node[above right=1pt] at (B) {$B(3,4)$};

% 动点 P(x,y) 选在角度 -45° 处: (2+1, 3-1) = (3, 2)
\coordinate (P) at (3, 2);
\fill (P) circle (1.5pt);
\node[below right=1pt] at (P) {$P(x,y)$};

% 向量 OP
\draw[->, line width=1.0pt] (0,0) -- (P);

% 变力 F: 垂直于 OP，指向第二象限 (-y, x) 方向，比例缩放
\coordinate (F_dir) at (-1.2, 1.8);
\draw[->, line width=1.1pt] (P) -- ++(-1.6, 2.4) node[above right=1pt] {$\vec{F}$};
